% Generated draft: requires mathematical and source review.
\documentclass[11pt]{article}
\usepackage[margin=25mm]{geometry}
\usepackage{fontspec}
\setmainfont{Latin Modern Roman}
\usepackage{amsmath,amssymb,mathtools,unicode-math}
\setmathfont{Latin Modern Math}
\usepackage{enumitem,needspace,xurl,hyperref}
\setlength{\parindent}{0pt}
\setlength{\parskip}{5pt}
\setlength{\emergencystretch}{3em}
\newcommand{\N}{\mathbb{N}}
\newcommand{\Z}{\mathbb{Z}}
\newcommand{\Q}{\mathbb{Q}}
\newcommand{\R}{\mathbb{R}}
\newcommand{\C}{\mathbb{C}}
\newcommand{\F}{\mathbb{F}}
\newcommand{\A}{\mathbb{A}}
\newcommand{\PP}{\mathbb{P}}
\newcommand{\E}{\mathbb{E}}
\newcommand{\eps}{\varepsilon}
\newcommand{\dd}{\,\mathrm{d}}
\begin{document}
\section*{UnsolvedMath ID: 20003311}
\textbf{Source problem code:} catalogue entry
\subsection{Definitions and mathematical statement}
Question 31, attributed to G. Margulis, asks that for every $\varepsilon>0$ the positive spectrum $[\varepsilon,\infty)\cap\{m(F)\}$ be finite. The local sentence does not define the forms $F$ or normalization of $m(F)$. The aid interprets it as the cubic multiplicative spectrum of unimodular lattices, not an assertion about arbitrary unnormalized functions.

In that interpretation take $B\in SL(3,\R)$, $L_B=B\Z^3$ and $F_B(z)=(Bz)_1(Bz)_2(Bz)_3$. Set
\[
 m(F_B)=\inf_{z\in\Z^3\setminus\{0\}}|F_B(z)|
 =\mu(L_B),\qquad
 \mu(L)=\inf_{0\ne v\in L}|v_1v_2v_3|.
\]
For an arbitrary invertible coefficient matrix, unimodular rescaling divides the product minimum by $|\det B|$; the original convention must be checked before using that extension. The reconstructed question is finiteness of the distinct values of $\mu(L)$ at least $\varepsilon$, not finiteness of the lattices themselves.

Let $A^+$ be the full positive determinant-one diagonal group and $\operatorname{sys}_\infty(L)=\inf_{0\ne v\in L}\|v\|_\infty$. The aid reports the exact identity
\[
 \mu(L)=\left(\inf_{a\in A^+}\operatorname{sys}_\infty(aL)\right)^3.
\]
With $K_r=\{L:\operatorname{sys}_\infty(L)\geq r\}$, its compact-threshold consequence is
\[
 \mathcal B_\varepsilon=\{L:\mu(L)\geq\varepsilon\}
 =\bigcap_{a\in A^+}a^{-1}K_{\varepsilon^{1/3}}.
\]
Mahler compactness and closedness make $\mathcal B_\varepsilon$ compact. This does not imply that its scalar minimum spectrum is finite: compactness alone permits infinitely many values.

The reported logical reduction separates orbit and spectrum information. If every bounded $A^+$ orbit is closed, every orbit in a positive threshold is compact. If, additionally, only finitely many compact closed orbits meet that threshold, the spectrum there is finite. Conversely, scalar-spectrum finiteness cannot alone give orbit finiteness without a finite-fiber statement for $\mu$ on the relevant orbit space. With such a finite-fiber hypothesis it yields finite orbit counts at each threshold; any further claimed closedness equivalence needs its own argument. No full equivalence is silently inferred from an invariant scalar.

For normalization the aid gives a Shanks cubic-order family of compact periodic orbits with exact minima
\[
 \mu(L_t)=\frac1{t^2+3t+9}.
\]
Its defining polynomial/order lattice and admissible integer range are not supplied. The summary also claims a closed threshold-count formula but does not print it. These remain reported family data and do not establish finiteness for all lattices.

The extraction then begins separate Section 5.4, defining, for $0\leq s\leq1$,
\[
 C_s=\{(\alpha,\beta)\in\R^2:
 \inf_{q\geq1}\max\{q^s\|q\alpha\|,
 q^{1-s}\|q\beta\|\}>0\},
\]
where $\|x\|$ is distance to the nearest integer. It identifies $C_{1/2}$ with simultaneous badly approximable vectors and ends with the fragment ``Since.'' This weighted-approximation context is not an added clause of Question 31, and its missing continuation is not reconstructed. The aid leaves the full positive-spectrum question open.
\subsection{Short English statement}
The aid translates one normalized reading into multiplicative lattice minima and proves compactness of each positive threshold. Finiteness of values still needs additional orbit information; a compact set alone does not give it. The explicit cubic-order family calibrates the minimum but does not settle the full spectrum.
\subsection{Sources}
\begin{itemize}
\item Local UnsolvedMath registry, permanent ID 20003311, source code catalogue entry; \texttt{aim-other-notes.json}, notes[117], source tag question.
\item AIM, \emph{Emerging applications of measure rigidity}, Question 31, G. Margulis, followed by Section 5.4: \url{https://aimath.org/WWN/measrigid/measrigid.pdf}.
\item Mahler compactness, Margulis bounded-orbit closedness and Shanks cubic orders are named by the supplied aid without complete bibliography or family specification.
\item Machine-generated research attempt 1: partial result, open at the research run, rigor labelled ``proved,'' low novelty confidence; registry partially solved. Not peer reviewed; independently verify before citation. Registry updated 14 August 2026.
\end{itemize}
\end{document}
